Although reference counting~\cite{classicRC} (RC) is one of the oldest
memory management techniques in computer science, it is not considered
a serious garbage collection technique in the functional programming
community, and there is plenty of evidence it is in general inferior
to tracing garbage collection algorithms. Indeed, high-performance
compilers such as ocamlopt and GHC use tracing garbage
collectors. Nonetheless, implementations of several popular
programming languages, e.g., Swift, Objective-C, Python, and Perl, use
reference counting as a memory management technique. Reference
counting is often praised for its simplicity, but many disadvantages
are frequently reported in the
literature~\cite{GCbook,GCsurvey}. First, incrementing and
decrementing reference counts every time a reference is created or
destroyed can significantly impact performance because they not only
take time but also affect cache performance, especially in a
multi-threaded program~\cite{biasedRC}. Second, reference counting
cannot collect circular~\cite{cyclesRC} or self-referential
structures. Finally, in most reference counting implementations, pause
times are deterministic but may still be unbounded~\cite{BoehmCostRC}.

In this paper, we investigate whether reference counting is a
competitive memory management technique for purely functional
languages, and explore optimizations for reusing memory, performing
destructive updates, and for minimizing the number of reference count
increments and decrements. The former optimizations in particular are
beneficial for purely functional languages that otherwise can only
perform functional updates. When performing functional updates,
objects often die just before the creation of an object of the same
kind.  We observe a similar phenomenon when we insert a new element
into a pure functional data structure such as binary trees, when we
use \emph{map} to apply a given function to the elements of a list or
tree, when a compiler applies optimizations by transforming abstract
syntax trees, or when a proof assistant rewrites formulas. We call it the
\emph{resurrection hypothesis}: many objects die just before the
creation of an object of the same kind. Our new optimization takes
advantage of this hypothesis, and enables pure code to perform
destructive updates in all scenarios described above when objects are
not shared.  We implemented all the ideas reported here in the new
runtime and compiler for the Lean programming language~\cite{Lean}. We
also report preliminary experimental results that demonstrate our new
compiler produces competitive code that often outperforms the code
generated by high-performance compilers such as ocamlopt and GHC
(\cref{sec:evaluation}).

Lean implements a version of the Calculus of Inductive
Constructions~\cite{coquand:huet:88,coquand:paulin:mohring:90}, and it
has mainly been used as a proof assistant so far. Lean has a
metaprogramming framework for writing proof and code automation, where
users can extend Lean using Lean itself~\cite{ebner17meta}. Improving
the performance of Lean metaprograms was the primary motivation for the work
reported here, but one can apply the techniques reported here
to general-purpose functional programming languages.
%% In the current Lean release, metaprograms are compiled
%% into bytecode and evaluated in a virtual machine. This virtual machine
%% uses a simple memory management technique based on reference
%% counting. Lean users have written many nontrivial metaprograms such as
%% a ring theory solver~\cite{RingSolver}, a decision procedure for
%% Presburger arithmetic~\cite{Presburger}, and a Fourier-Motzkin solver
%% for real linear arithmetic, to cite a few. Using the current stable
%% version of Lean, these useful metaprograms are not competitive with
%% similar proof automation implemented in functional programming
%% languages with an efficient compiler such as Haskell and OCaml. The
%% primary source of inefficiency in Lean metaprograms is the virtual
%% machine interpretation overhead and the naive reference counting
%% approach used.
%% Although the efficiency of Lean metaprograms was the
%% primary motivation for our work, we believe the techniques reported
%% here are applicable to general-purpose programming.

We describe our approach as a series of refinements starting from
\pureIR{}, a simple intermediate representation for eager and purely
functional languages (\cref{sec:pure}). We remark that in Lean and
\pureIR{}, it is not possible to create cyclic data structures. Thus,
one of the main criticisms against reference counting does not
apply. From \pureIR{}, we obtain \rcIR{} by adding explicit instructions
for incrementing ($\kw{inc}$) and decrementing ($\kw{dec}$) reference
counts, and reusing memory (\cref{sec:rc}). The inspiration for explicit
RC instructions comes from the Swift compiler, as does the notion
of \emph{borrowed} references. In contrast to standard
(or \emph{owned}) references, of which there should be exactly as many
as the object's reference counter implies, borrowed references do not
update the reference counter but are \emph{assumed} to be kept alive
by a surrounding owned reference, further minimizing the number of
$\kw{inc}$ and $\kw{dec}$ instructions in generated code.

We present a simple compiler from \pureIR{} to \rcIR{},
discussing heuristics for inserting destructive updates, borrow
annotations, and \kw{inc}/\kw{dec} instructions (\cref{sec:compiler}).
Finally, we show that our approach is compatible with existing techniques
for performing destructive updates on array and string values, and
propose a simple and efficient approach for thread-safe reference
counting (\cref{sec:runtime}).
%% (Leo): the following information is not relevant for most readers
%% This is important because, in the Lean
%% system, we compile function definitions and theorems in parallel using
%% multi-threaded tasks. Complex theorems make heavy use of proof
%% automation, and may take a long time to compile. Multi-threaded
%% compilation significantly reduces the compilation time.

\emph{Contributions}. We present a reference counting scheme optimized for and
used by the next version of the Lean programming language.
\begin{itemize}
\item We describe how to reuse allocations in both user code and
  language primitives, and give a formal reference-counting semantics that
  can express this reuse.
\item We describe the optimization of using borrowed references.
  %% , and how to
  %% formalize them as intuitionistic variables in an otherwise linear type system.
\item We define a compiler that implements all these steps. The compiler is implemented
  in Lean itself and the source code is available.
\item We give a simple but effective scheme for avoiding atomic reference count
  updates in multi-threaded programs.
\item We compare the new Lean compiler incorporating these ideas with other
  compilers for functional languages and show its competitiveness.
\end{itemize}

\section{Examples}
\label{sec:examples}

In reference counting, each heap-allocated value contains a reference
count. We view this counter as a collection of tokens. The \kw{inc}
instruction creates a new token and \kw{dec} consumes it. When a function
takes an argument as an owned reference, it is responsible for
consuming one of its tokens. The function may consume the owned
reference not only by using the \kw{dec} instruction, but also by storing
it in a newly allocated heap value, returning it, or passing it to
another function that takes an owned reference. We illustrate our intermediate
representation (IR) and the use of owned and borrowed references with a series of
small examples.

The identity function $\textit{id}$ does not require any RC operation
when it takes its argument as an owned reference.
\begin{lstlisting}
id x = ret x
\end{lstlisting}
%% Now, consider the function $\textit{mkTriple}$ that takes three arguments $x$, $y$, and $z$,
%% and returns $(x, (y, z))$. Again, $\textit{mkTriple}$ does not require any RC operation
%% when the arguments are owned references.
%% \begin{lstlisting}
%% mkTriple x y z = let p = Pair y z; let q = Pair x p; ret q
%% \end{lstlisting}
%% Here, $y$ and $z$ are consumed by the constructor $\textit{Pair}\ y\ z$ which allocates a new pair in the heap.
%% Similarly, $x$ and $q$ are consumed by $\textit{Pair}\ x\ p$, and $q$ by $\kw{ret}\ q$.
As another example, consider the function $\textit{mkPairOf}$ that takes $x$ and returns the pair $(x, x)$.
\begin{lstlisting}
mkPairOf x = inc x; let p = Pair x x; ret p
\end{lstlisting}
It requires an \kw{inc} instruction because two tokens for $x$ are consumed (we will also say that ``$x$ is consumed'' twice).
The function $\textit{fst}$ takes two arguments $x$ and $y$, and
returns $x$, and uses a $\kw{dec}$ instruction for consuming the unused $y$.
\begin{lstlisting}
fst x y = dec y; ret x
\end{lstlisting}
The examples above suggest that we do not need any RC operation when
we take arguments as owned references and consume them exactly
once. Now we contrast that with a function that only inspects its argument:
the function \textit{isNil} \textit{xs} returns
true if the list \textit{xs} is empty and false otherwise. If the argument $\textit{xs}$ is taken as an
owned reference, our compiler generates the following code
\begin{lstlisting}
  isNil xs = case xs of
    (Nil -> dec xs; ret true)
    (Cons -> dec xs; ret false)
\end{lstlisting}
We need the \kw{dec} instructions because a function must consume all
arguments taken as owned references.  One may notice that decrementing
$\text{xs}$ immediately after we inspect its constructor tag is
wasteful. Now assume that instead of taking the ownership of an RC
token, we could borrow it from the caller.  Then, the callee would not
need to consume the token using an explicit \kw{dec} operation. Moreover,
the caller would be responsible for keeping the borrowed value alive.
This is the essence of \emph{borrowed references}: a borrowed reference does not actually keep the referenced value
alive, but instead asserts that the value is kept alive by another, owned reference.
Thus, when \textit{xs} is a borrowed reference, we compile \textit{isNil} into our IR as
\begin{lstlisting}
  isNil xs = case xs of (Nil -> ret true) (Cons -> ret false)
\end{lstlisting}

As a less trivial example, we now consider the function \textit{hasNone} \textit{xs} that,
given a list of optional values, returns \emph{true} if \textit{xs} contains a \textit{None} value.
This function is often defined in a functional language as
\begin{lstlisting}
  hasNone []            = false
  hasNone (None : xs)   = true
  hasNone (Some x : xs) = hasNone xs
\end{lstlisting}
Similarly to \textit{isNil}, \textit{hasNone} only inspects its argument. Thus if \textit{xs} is taken
as a borrowed reference, our compiler produces the following RC-free IR code for it
\begin{lstlisting}
  hasNone xs = case xs of
    (Nil -> ret false)
    (Cons -> let h = proj$_{\textit{head}}$ xs; case h of
      (None -> ret true)
      (Some -> let t = proj$_{\textit{tail}}$ xs; let r = hasNone t; ret r))
\end{lstlisting}
Note that our \kw{case} operation does not introduce binders. Instead, we use explicit instructions $\kw{proj}_i$ for
accessing the head and tail of the \textit{Cons} cell. We use suggestive names
for cases and fields in these initial examples, but will later use indices instead.
Our borrowed inference heuristic discussed in \cref{sec:compiler} correctly tags $\textit{xs}$ as a borrowed parameter.

When using owned references, we know at run time whether a
value is shared or not simply by checking its reference counter.
We observed we could leverage this information and minimize
the amount of allocated and freed memory for constructor values such as a
list \textit{Cons} value. Thus, we have added two additional instructions to our IR:
$\kw{let}\ y = \kw{reset}\ x$ and $\kw{let}\ z = (\kw{reuse}\ y\ \kw{in}\ \kw{ctor}_i\ \many{w})$.
The two instructions are used together; if $x$ is a shared value, then $y$ is set to
a special reference $\nullptr$,
and the $\kw{reuse}$ instruction just allocates a new constructor value $\kw{ctor}_i\ \many{w}$.
If $x$ is not shared, then $\kw{reset}$ decrements the
reference counters of the components of $x$, and $y$ is set to $x$. Then, $\kw{reuse}$ reuses
the memory cell used by $x$ to store the constructor value $\kw{ctor}_i\ \many{w}$.
We illustrate these two instructions with the IR code for the list \textit{map} function
generated by our compiler as shown in \cref{sec:compiler}. The code uses our
actual, positional encoding of cases, constructors, and fields as described in
the next section.
\begin{lstlisting}
  map f xs = case xs of
    (ret xs)
    (let x = proj_1 xs; inc x; let s = proj_2 xs; inc s;
     let w = reset xs;
     let y = f x; let ys = map f s;
     let r = (reuse w in ctor_2 y ys); ret r)
\end{lstlisting}
%% If \textit{xs} is \textit{Nil}, then \textit{map} just returns \textit{xs}.
%% Otherwise, it extracts the head \textit{x} and tail \textit{s} from \textit{xs} using
%% the projection $\kw{proj}_i$, and uses the \inc instruction to create owned references to
%% these two values. Then, $\textit{reset}\ xs$ consumes \textit{xs} and if it
We remark that if the list referenced by $\textit{xs}$ is not shared, the code above
does not allocate any memory. Moreover, if \textit{xs} is a nonshared list of list of integers,
then $\textit{map}\ (\textit{map}\ \textit{inc})\ \textit{xs}$ will not allocate
any memory either.
This example also demonstrates it is not a good idea, in general, to fuse
$\kw{reset}$ and $\kw{reuse}$ into a single instruction: if we removed the
$\kw{let}\ w = \kw{reset}\ \textit{xs}$ instruction and directly used
\textit{xs} in \kw{reuse},
then when we execute the recursive application $\textit{map}\ f\ s$, the reference counter for $s$ would
be greater than $1$ even if the reference counter for $\textit{xs}$ was 1.
We would have a reference from $\textit{xs}$ and another from $s$, and memory reuse would not occur in the recursive applications.
Note that removing the $\kw{inc}\ s$ instruction is incorrect when
$\textit{xs}$ is a shared value.
Although the $\kw{reset}$ and $\kw{reuse}$ instructions can in general be used for reusing memory
between two otherwise unrelated values, in examples like \textit{map} where the
reused value has a close semantic connection to the reusing value, we will use
common functional vocabulary and say that the list is being \emph{destructively updated} (up to the first shared cell).

As another example, a zipper is a technique for traversing and efficiently updating data structures, and it is particularly useful for
purely functional languages. For example, the list zipper is a pair of lists, and it allows one to
move forward and backward, and to update the current position.
The \textit{goForward} function is often defined as
\begin{lstlisting}
   goForward ([], bs)   = ([], bs)
   goForward (x : xs, bs) = (xs, x : bs)
\end{lstlisting}
In most functional programming languages, the second equation allocates a
new pair and \textit{Cons} value.  %% Moreover, if the input argument is
%% not shared, the input pair becomes dead after the function returns.
The functions \textit{map} and \textit{goForward} both satisfy our
resurrection hypothesis. Moreover, the result of a \textit{goForward}
application is often fed into another \textit{goForward} or
\textit{goBackward} application.  Even if the initial value was shared,
every subsequent application takes a nonshared pair, and
memory allocations are avoided by the code produced by our compiler.
%% Our compiler produces the following IR code for \textit{goForward}
\begin{lstlisting}
   goForward p = case p of
     (let xs = proj_1 p; inc xs;
      case xs of
        (ret p)
        (let bs = proj_2 p; inc bs;
         let c_1 = reset p;
         let x = proj_1 xs; inc x; xs' = proj_2 xs; inc xs';
         let c_2 = reset xs;
         let bs' = (reuse c_2 in ctor_2 x bs);
         let r   = (reuse c_1 in ctor_1 xs' bs'); ret r))
\end{lstlisting}

%% Although reference counting is one of the oldest memory
%% management techniques in computer science, it is not considered a
%% serious garbage collection technique in the functional programming
%% community, and there is plenty of evidence that it is inferior to
%% tracing garbage collecting algorithms. Moreover, reference counting is not used in
%% high-performance compilers such as ocamlopt and GHC.

%% is often praised for its simplicity, but many disadvantages
%% are frequently reported in the literature. First, incrementing and
%% decrementing reference counts every time a reference is created or
%% destroyed can significantly impact performance because they not only
%% take time but also affect cache performance. Second, reference
%% counting cannot collect circular or self-referential
%% structures. Finally, in most reference counting implementations,
%% unbounded pause times may occur.

%% Reference counting is expensive - every time you manipulate pointers to an object you need to update and check the reference count. Pointer manipulation is frequent, so this slows your program and bloats the code size of compiled code.

%% They cannot collect so-called circular, or self-referential structures.


%% OCaml garbage collector is a modern hybrid generational/incremental collector which outperforms hand-allocation in most cases.

%% GHC uses a generational garbage collector and so the short lived garbage is quickly and cheaply collected.

%% Reference counting is considered a bad solution:

%% Incrementing and decrementing reference counts every time a reference is created or destroyed can significantly impede performance.
%% Not only do the operations take time, but they affect cache performance.

%% Even read-only operations like calculating the length of a list
%% require a large number of reads and writes for reference updates with
%% naive reference counting.

%% It is used in C++, Objective-C, Perl, Tcl.

%% poor cache performance


%% Lean implements a version of the Calculus
%% of Inductive Constructions

%% Lean also has parallel compilation and checking of proofs,
%% and provides a server mode that supports a continuous compilation and rich user interaction in
%% editing environments such as Emacs and Visual Studio Code.

%% Here we describe the metaprogramming framework currently
%% in use in Lean

%% Other systems use meta programming.

%% There are a number of tactic languages currently on offer, including Ltac
%% [Delahaye 2000, 2002], Mtac [Ziliani et al. 2015], and MetaCoq [Ziliani et al. 2017] for Coq and
%% Eisbach [Matichuk et al. 2016] for Isabelle.

%% For effective metaprogramming, it is crucial to have efficient means of constructing objectlanguage
%% expressions in the metaprogramming language. Our language incorporates mechanisms
%% for direct reflection [Barzilay 2006], in the sense that the same syntax and elaboration mechanisms
%% used for the object language are available in the metaprogramming language. In this sense, it is
%% an advance over the reflection mechanism adopted in Agda [van der Walt and Swierstra 2012] for
%% similar purposes. (The current version of Agda, 2.5.2, h


%% Text of paper \ldots

%% \begin{itemize}
%% \item context: eager+pure: constraints \& opportunities
%% \item basic runtime model: refcounter, allocation, deallocation
%% \item destructive updates
%% \item explicit RC instructions
%% \item concurrency
%% \item evaluation
%% \end{itemize}

%%% Local Variables:
%%% mode: latex
%%% TeX-master: "refcount"
%%% End:
